\scanpage{153}
\mspage{98}
que $\rho$ et $i$ soient adjoints, \deleted{Mais le cas de $\mathrm{Cat}$ \inserted{(si je ne m'y suis \unsure{trompé})} \unsure{nous montre} qu'ils peuvent être adjoints} $i$ étant l'adjoint à droite de $\rho$ :
\diagram{\begin{tikzcd}[ampersand replacement=\&]
\mathrm{Ho}^{\mathrm{lc}}_W(S)\arrow[r,bend right=12,"i"'] \& \mathrm{Ho}^{\mathcal L}_W(S)\arrow[l,bend right=12,"\rho"']\text{,}
\end{tikzcd}}
i.e on cherche
\[
\begin{aligned}
i\rho &\xleftarrow{\overset{\text{\uline{pas} iso}}{\unsure{i|\widetilde{\imath}}}}\mathrm{id} &&\text{dans }\mathrm{Ho}^{\mathcal L}_W &&(\text{loc. par }W_S^{\mathcal L})\\
\rho i &\xrightarrow[p]{\mathrm{iso}}\mathrm{id} &&\text{dans }\mathrm{Ho}^{\mathrm{lc}}_W &&(\text{loc par }W_S)
\end{aligned}
\]
donc pour $X$ dans $\mathcal M_{/S}$, on cherche
\[
\begin{aligned}
\unsure{i_X|\widetilde{\imath}_X}:&\quad X\xrightarrow{\deleted{\unsure{\rho i}}} C(X) &&\substack{\text{tout au moins}\\\text{dans }(\mathcal M_{/S})(W_S^{\mathcal L})^{-1}}\\
p_X:&\quad X\xleftarrow{\deleted{\unsure{p}}} C(X) &&\substack{\text{tout au moins dans}\\(\mathcal M_{/S})(W_S)^{-1}}\text{.}
\end{aligned}
\]
On choisira pour $i_X$ la flèche ``canonique'' de $X$ vers $C(X)$, définie dans $\mathcal M_{/S}$ lui-$\eqsim$, sans avoir à localiser. Mais dans le ``cas Quillen'' (quand $C$ n'est pas donné comme foncteur $\mathcal M_{/S}\longrightarrow\underline{\mathrm{Fib}}\,S$), il faut voir que cet $i_X$ est fonctoriel en $X$, quand on le regarde dans le localisé $(\mathcal M_{/S})(W_S)^{-1}$. Mais c'est contenu dans la construction de $C$ comme \uline{foncteur}
\[
\mathcal M_{/S}\longrightarrow(\underline{\mathrm{Fib}}\,S)(W_S)^{-1}
\]
On notera que $i_X$ est dans $W_S$ sans plus, pas dans $W_S^{\mathcal L}$, donc ne définit pas en général un iso de $\mathrm{Ho}^{\mathcal L}_W$.

On définit $p_X$ à l'aide de $i_X$, comme l'inverse de $i_X$ \uline{dans} $(\mathcal M_{/S})(W_S)^{-1}$, qui existe car
% archive sheet number: 152
